\chapter{Git と GitHub の基本}
\label{app:git}
\index{git@Git|textbf}
\index{github@GitHub|textbf}

\section{リポジトリ・ステージング・コミットの関係}

Git は、手元のファイルの変更履歴を記録するためのソフトウェアである。これに対して GitHub は、Git リポジトリを置いて共有するためのサービスである。

\begin{itemize}
\item Git: ローカルマシンでも単独で使えるバージョン管理システム
\item GitHub: Git リポジトリをホストし、共有、Issue、Pull Request などを提供するサービス
\end{itemize}

したがって、\code{git init}、\code{git add}、\code{git commit} は Git の話であり、GitHub 上のリポジトリを \code{clone} で取得したり、\code{push} で送ったりする段階で GitHub が関係する。

\subsection{リポジトリと公開範囲}
\index{りぽじとり@リポジトリ|textbf}
\index{こうかいはんい@公開範囲|textbf}

リポジトリは、コード、データ、履歴をまとめて管理する単位である。GitHub へ置く場合は、公開範囲も意識する必要がある。

\begin{itemize}
\item 公開（\code{public}）: 誰でも閲覧できる
\item 非公開（\code{private}）: 明示的に許可された利用者だけが閲覧できる
\item 内部（\code{internal}）: 主に GitHub Enterprise で、同じ Enterprise に属するメンバーへ公開する。単一の Organization 内だけに限定する設定とは異なる
\end{itemize}

授業用の配布物、研究中のコード、機密データを含むリポジトリでは、公開範囲の誤設定が情報漏えいにつながる。GitHub で新しいリポジトリを作成する際には、公開範囲（visibility）の確認が必要である。

\subsection{作業ツリー・ステージングエリア・コミット履歴}
\index{すてーじんぐえりあ@ステージングエリア|textbf}
\index{わーきんぐつりー@ワーキングツリー|textbf}

Git では、作業中のファイルと履歴へ記録する内容を分けて管理する。変更を履歴へ記録する手順は、次の 3 段階からなる。

\begin{enumerate}
\item 作業ファイルの編集
\item \code{git add} による記録対象の選択
\item \code{git commit} による履歴の確定
\end{enumerate}

この 2 段階目にある、次のコミットへ記録する内容を保持する領域をステージングエリアという。\code{git add} は作業ツリーの内容をステージングエリアへ登録し、\code{git commit} は登録済みの内容を履歴へ記録する。

Git が管理する状態は、次の 3 層に分けられる。

\begin{itemize}
\item 作業ツリー（working tree）: エディタで直接編集しているファイル
\item ステージングエリア（staging area、index）: 次のコミットへ入れる変更を一時的に集めた領域
\item コミット履歴（commit history）: 確定済みの変更が記録された履歴
\end{itemize}

たとえば、ファイルを編集した直後の変更は作業ツリーだけに存在する。\code{git add} を実行すると、その時点の内容がステージングエリアへ登録される。\code{git commit} は、ステージングエリアの内容を履歴として確定する操作である。\code{git add} の後に同じファイルを編集した場合、追加分は作業ツリーだけに残り、ステージ済みの内容には含まれない。

\section{ローカル履歴の作成}

\subsection{コミットへ記録する名前とメールアドレスの設定}
\index{git config@\code{git config}|textbf}

コミットへ記録する作成者名とメールアドレスを設定する。

\begin{lstlisting}[numbers=none, xleftmargin=2\zw]
$ git config --global user.name "Your Name"
$ git config --global user.email "your_email@example.com"
\end{lstlisting}

これをしておかないと、コミット時に誰の変更なのかが分からなくなる。

\subsection{新しいリポジトリの作成}
\index{git init@\code{git init}|textbf}

まだ Git で管理されていないディレクトリを初期化するコマンドが \code{git init} である。

\begin{lstlisting}[numbers=none, xleftmargin=2\zw]
$ mkdir analysis_work
$ cd analysis_work
$ git init
$ git status
\end{lstlisting}

\code{git status} は、未追跡、変更済み、コミット待ちのファイルを一覧できる。編集やコミットの前後で実行すると、各変更がどの段階にあるかを確認できる。

\subsection{編集、確認、記録の流れ}
\index{git status@\code{git status}|textbf}
\index{git diff@\code{git diff}|textbf}
\index{git add@\code{git add}|textbf}
\index{git commit@\code{git commit}|textbf}

日常的な流れは、編集して、差分を確認して、必要なものだけをコミットする、である。

\begin{lstlisting}[numbers=none, xleftmargin=2\zw]
$ git status
$ git diff
$ git add analyze.py
$ git diff --cached
$ git commit -m "Add first analysis script"
$ git log --oneline
\end{lstlisting}

\code{git diff} は、作業ツリーとステージングエリアの差分を表示する。\code{git diff --cached} は、ステージングエリアと直前のコミットの差分を表示する。\code{git log --oneline} は、コミット履歴を 1 件 1 行で表示する。

\subsubsection{\code{git status} の読み方}

\code{git status} の表示では、少なくとも次の 3 種類を区別する。

\begin{itemize}
\item 未追跡（untracked）: まだ Git が追跡していない新規ファイル
\item 変更済み（modified）: 追跡中だが、変更内容をまだ \code{add} していないファイル
\item コミット予定（changes to be committed）: \code{add} 済みで、次のコミットへ入る変更
\end{itemize}

\code{git status} の出力から、各ファイルが未追跡、変更済み、コミット予定のいずれに該当するかを判別できる。

\subsubsection{\code{git add} による変更内容のステージング}

\code{git add} は「変更を次のコミット候補として選ぶ」コマンドである。保存ではない。したがって、\code{git add} した後でも、コミット前にさらに編集すれば、その追加分はまだコミット対象へ入っていない。

\begin{lstlisting}[numbers=none, xleftmargin=2\zw]
$ git diff
$ git add analyze.py
$ git diff --cached
\end{lstlisting}

この 2 種類の \code{diff} を区別すると、作業ツリーと次のコミット候補の差を把握できる。

\subsection{\code{.gitignore} による除外設定}
\index{gitignore@\code{.gitignore}|textbf}

生成物まで毎回コミットすると、履歴が見にくくなる。Python 解析なら、仮想環境、キャッシュ、ビルド成果物などは除外することが多い。

\subsubsection{\code{.gitignore} の記述}

\code{.gitignore} は、Git に「このパターンのファイルは追跡しなくてよい」と教えるためのファイルである。リポジトリのルートに置くことが多い。Python 解析用の最小例は次のようになる。

\begin{lstlisting}[numbers=none, xleftmargin=2\zw, title={.gitignore}]
.venv/
__pycache__/
*.pyc
build/
dist/
.pytest_cache/
\end{lstlisting}

このファイルを作ると、まだ追跡されていない \code{.venv/} や \code{build/} は無視され、通常の \code{git status} には表示されない。すでに追跡中のファイルには、この設定は適用されない。

\subsubsection{すでに追跡済みなら \code{git rm --cached} が必要}
\index{git rm cached@\code{git rm --cached}|textbf}

\code{.gitignore} は「これから追跡しない」ための設定である。すでに一度コミットしたファイルには、そのままでは効かない。

\begin{lstlisting}[numbers=none, xleftmargin=2\zw]
$ git rm --cached -r .venv __pycache__ build
$ git commit -m "Stop tracking generated files"
\end{lstlisting}

\code{git rm --cached} は、作業ディレクトリの実ファイルを残したまま、インデックスに登録されたパスだけを削除する。実行後の削除をコミットすると、そのパスは Git の追跡対象から外れる。

\subsubsection{\code{uv} プロジェクトにおけるコミット対象}

\code{uv} を使うプロジェクトでは、通常 \code{.venv/} はコミットしない。依存関係を記述する \code{pyproject.toml} と除外規則を記述する \code{.gitignore} はコミットする。アプリケーションや再現可能な解析環境では、解決済みの依存関係を記録する \code{uv.lock} も通常はコミットする。ライブラリでは、リポジトリ内のテストや開発環境をロックファイルで再現するかどうかを運用方針として定める。\code{.python-version} は、プロジェクトで使用する Python の版を共有する場合にコミットする。

\section{GitHub への接続}

\subsection{既存リポジトリの取得}
\index{git clone@\code{git clone}|textbf}

\code{git clone} は、GitHub 上のプロジェクトや教材を履歴ごと手元へ複製するコマンドである。

\begin{lstlisting}[numbers=none, xleftmargin=2\zw]
$ git clone git@github.com:example/project.git
$ cd project
$ git remote -v
\end{lstlisting}

\code{git clone} は、リポジトリのファイルと履歴をローカルへ複製する。\code{git remote -v} は、接続先として登録された URL を表示する。本書の例では、取得と push の双方に SSH URL を用いる。公開鍵の準備は付録~\ref{app:ssh} で説明する。

\subsection{ローカルリポジトリの GitHub への初回公開}
\index{りもーとりぽじとり@リモートリポジトリ|textbf}
\index{git remote@\code{git remote}|textbf}

既存リポジトリの取得とは逆に、ローカルで作成した履歴を新しい GitHub リポジトリへ送る場合がある。まず GitHub 上に空のリポジトリを作成し、次にローカルリポジトリへ接続先を登録して push する。

\subsubsection{GitHub 上の空リポジトリの作成}

GitHub の Web 画面で \code{New repository} を開き、リポジトリ名と公開範囲を決める。ローカルにすでに履歴がある場合は、README、\code{.gitignore}、ライセンスを GitHub 側で追加せず、コミットを持たない空のリポジトリを作成する。

README や \code{.gitignore} を GitHub 側で追加すると、リモート側にローカルとは別のコミットが作られる。その場合、双方の履歴を統合してから push する必要がある。空のリモートリポジトリであれば、この統合は不要である。

\subsubsection{README や LICENSE を先に作っていた場合}

GitHub 側で README や LICENSE を作成した場合、リモートリポジトリにはローカルと異なるコミットが存在する。リモート側の履歴をローカルへ取り込んで統合した後に push する。

\begin{lstlisting}[numbers=none, xleftmargin=2\zw]
$ git branch -M main
$ git remote add origin git@github.com:owner/project.git
$ git fetch origin
$ git merge --allow-unrelated-histories origin/main
$ git push -u origin main
\end{lstlisting}

\code{--allow-unrelated-histories} が必要なのは、手元のリポジトリと GitHub 側のリポジトリが別々に作られており、履歴の出発点が一致していないからである。この \code{fetch} と \code{merge} によって、GitHub 側の README や LICENSE のコミットを自分の履歴へ取り込んでから、その続きとして push できるようになる。

ローカルとリモートの双方で同じ名前の README や LICENSE を作成している場合、\code{git merge} によって競合が生じることがある。競合箇所を編集し、\code{git add} と \code{git commit} で統合結果を記録してから、\code{git push -u origin main} を実行する。

\subsubsection{接続先の登録と最初の push}

GitHub が表示する URL を確認し、ローカルリポジトリへ \code{origin} という名前で登録する。既定ブランチを \code{main} とする例を次に示す。

\begin{lstlisting}[numbers=none, xleftmargin=2\zw]
$ git status
$ git branch --show-current
$ git branch -M main
$ git remote add origin git@github.com:owner/project.git
$ git remote -v
$ git push -u origin main
\end{lstlisting}

\code{git branch --show-current} は現在のブランチ名を表示する。\code{git branch -M main} は現在のブランチ名を \code{main} へ変更する。既存のブランチ名を維持する場合は、以後のコマンドにその名前を指定する。

\code{git remote add origin ...} は、GitHub の URL を \code{origin} という名前で登録する。\code{git remote -v} は登録後の取得用 URL と送信用 URL を表示する。\code{git push -u origin main} の \code{-u} は、ローカルの \code{main} に対する上流ブランチを \code{origin/main} に設定する。この設定後は、接続先とブランチ名を省略した \code{git push} と \code{git pull} を利用できる。

\subsubsection{\code{origin} がすでにある場合}

すでに \code{origin} が登録されている状態で \code{git remote add origin ...} を実行すると、エラーになる。その場合は、いま何が設定されているかを確認してから、必要なら URL を差し替える。

\begin{lstlisting}[numbers=none, xleftmargin=2\zw]
$ git remote -v
$ git remote set-url origin git@github.com:owner/project.git
$ git remote -v
\end{lstlisting}

教材を別のリポジトリから複製して流用した場合や、HTTPS から SSH へ切り替えたい場合に、この操作が必要になることがある。

\subsection{GitHub との同期}
\index{git pull@\code{git pull}|textbf}
\index{git push@\code{git push}|textbf}

ローカルでコミットしただけでは、まだ GitHub には反映されない。GitHub 側から最新を取り込むのが \code{pull}、手元のコミットを送るのが \code{push} である。

\begin{lstlisting}[numbers=none, xleftmargin=2\zw]
$ git pull --ff-only
$ git push origin main
\end{lstlisting}

\code{git pull --ff-only} は、ローカルブランチをリモートブランチまで早送りできる場合にだけ更新する。ローカルとリモートの両方に異なるコミットがあり履歴が分岐している場合は停止するため、意図しないマージコミットを作らない。

\subsection{GitHub で利用する HTTPS URL と SSH URL}
\index{github@GitHub!https@HTTPS}
\index{github@GitHub!ssh@SSH}

GitHub のリポジトリ URL には、少なくとも HTTPS と SSH の 2 通りがある。

\begin{lstlisting}[numbers=none, xleftmargin=2\zw]
https://github.com/owner/repo.git
git@github.com:owner/repo.git
\end{lstlisting}

HTTPS URL は HTTPS 認証、SSH URL は SSH 鍵認証を用いる。一つのリポジトリでは、登録したリモート URL と認証情報の対応を記録する。
本書では、\code{git clone}、\code{pull}、\code{push} に用いる認証方式を SSH に統一する。公開リポジトリの取得には HTTPS URL も利用できるが、push や非公開リポジトリの取得には認証が必要になる。SSH URL に統一すれば、操作ごとに認証方式を切り替える必要がない。

\section{認証と最小ワークフロー}

\subsection{GitHub の認証: SSH 鍵を基本にし、必要なら PAT}
\index{github@GitHub!にんしょう@認証|textbf}
\index{pat@個人用アクセストークン（PAT）|textbf}

GitHub では、HTTPS 経由の Git 操作にアカウントのパスワードを使用できない。SSH URL では SSH 鍵、HTTPS URL では個人用アクセストークン（PAT）または認証情報ヘルパーを用いる\footnote{\href{https://docs.github.com/en/authentication/keeping-your-account-and-data-secure/about-authentication-to-github}{GitHub Docs, ``About authentication to GitHub''}（2026 年 8 月 25 日閲覧）。}。本書では SSH 鍵を標準とし、HTTPS と PAT の組み合わせを代替手段として扱う。

\subsubsection{SSH 鍵による認証}

SSH を使う場合は、付録~\ref{app:ssh} の手順で公開鍵を用意し、それを GitHub の \code{SSH and GPG keys} へ登録する。GitHub CLI の \code{gh auth login --git-protocol ssh --web} は、認証と公開鍵の登録に必要な操作を対話形式で案内する。

\subsubsection{PAT による認証}

SSH 鍵を使えない事情があり、HTTPS で GitHub を使う場合は、プライベートリポジトリへの push や pull で PAT を使う。

\begin{enumerate}
\item GitHub の \code{Settings} への移動
\item \code{Developer settings} から \code{Personal access tokens} への移動
\item Fine-grained token または classic token の作成
\item 対象リポジトリと必要最小限の権限の選択
\item 発行済みトークンの安全な保管
\end{enumerate}

認証時には、ユーザー名として GitHub のアカウント名を入力し、パスワードを求められた欄へ PAT を入力する。組織がトークンの利用制限や承認を設定している場合には、その方針にも従う必要がある。

Fine-grained PAT では、対象となるリポジトリと権限を個別に限定できる。GitHub は、可能であれば personal access token (classic) よりも fine-grained PAT を推奨している\footnote{\href{https://docs.github.com/en/authentication/keeping-your-account-and-data-secure/managing-your-personal-access-tokens}{GitHub Docs, ``Managing your personal access tokens''}（2026 年 8 月 25 日閲覧）。}。ただし、組織の方針によって管理者の承認が必要な場合があり、利用する機能が fine-grained PAT に対応しているかも確認する。

SAML SSO を利用する組織では、SSH 鍵や PAT に組織リソースへのアクセス許可が必要になる場合がある。Personal access token (classic) は作成後に SSO への承認が必要であり、fine-grained PAT は作成時に組織を対象として指定する\footnote{\href{https://docs.github.com/authentication/authenticating-with-saml-single-sign-on/authorizing-a-personal-access-token-for-use-with-saml-single-sign-on}{GitHub Docs, ``Authorizing a personal access token for use with single sign-on''}（2026 年 8 月 25 日閲覧）。}。実際の承認手順と利用可能な認証方式は、対象組織の方針を確認する。

\subsection{変更確認から同期までの最小ワークフロー}

変更の確認から GitHub との同期までは、次の順序で行う。

\begin{enumerate}
\item \code{git status} による状況の確認
\item \code{git diff} による変更内容の確認
\item \code{git add} による記録対象の選択
\item \code{git diff --cached} によるコミット内容の再確認
\item \code{git commit -m "..."} による履歴への記録
\item GitHub との同期に用いる \code{git pull --ff-only} と \code{git push}
\end{enumerate}

この流れは、教材の取得、\code{uv} プロジェクトの再現、共同研究でのコード共有に共通する。\code{git add} と \code{git commit} の前後で \code{status} と \code{diff} を確認し、記録対象と実際の差分を一致させる。

公開範囲やアクセス条件の詳細は、GitHub のリポジトリに関する公式資料~\cite{github_visibility} を参照する。
